\section{Problem-Solving Agents: el caso de MLAgentBench}

\subsection{Convertir «saber conversar» en «saber hacer»}
El expositor presentó una dirección representativa de los problem-solving agents: completar un ciclo de extremo a extremo en tareas reales de aprendizaje automático, que incluye leer la tarea, modificar el código, ejecutar experimentos, analizar registros e iterar la estrategia.

\begin{figure}[H]
\centering
\includegraphics[width=0.93\textwidth]{frame-06-groupchat.jpg}
\caption{Planeación y resolución de tareas complejas: de la conversación a un pipeline ejecutable}
\end{figure}
\noindent\footnotesize Evidencia en pantalla: los subtítulos se ubican aproximadamente en 00:10:02--00:13:10, donde el expositor discute la resolución de tareas del benchmark, la ejecución en bash y la evaluación automática.\normalsize

\begin{importantbox}{Valor del benchmark: convertir la «discusión sobre capacidades» en un «experimento reproducible»}
Sin un benchmark, la mejora del agente suele quedarse en la narrativa de una demostración.  
Con un benchmark, es posible comparar la ganancia real entre distintos modelos, distintas estrategias y distintas cadenas de herramientas, y descubrir la distribución de los fallos.
\end{importantbox}

\subsection{Objetivos de evaluación y descomposición de fallos}
\begin{table}[H]
\centering
\begin{tabular}{p{2.8cm}p{4.0cm}p{4.8cm}}
\toprule
\textbf{Dimensión de evaluación} & \textbf{Señal observable} & \textbf{Falla típica} \\
\midrule
Grado de cumplimiento de la tarea & Si la métrica alcanza el objetivo, si el envío es válido & Desviación del objetivo, malentendido de la tarea \\
Confiabilidad de la ejecución & Tasa de éxito de los comandos, número de reintentos, tasa de tiempos de espera agotados & Errores de dependencia del entorno, fallas en la llamada a herramientas \\
Calidad de la estrategia & Eficiencia de la iteración, cobertura de la búsqueda, eficacia de la reflexión & Prueba y error a ciegas, óptimo local \\
Seguridad y restricciones & Comportamiento fuera de los permisos otorgados, efectos secundarios anómalos, capacidad de auditoría & Operación no autorizada, imposibilidad de revertir \\
\bottomrule
\end{tabular}
\caption{La evaluación de un problem-solving agent no solo mira la «puntuación final», sino también la calidad de la trayectoria de ejecución}
\end{table}

\begin{knowledgebox}{Por qué las cyber tasks se convierten en un escenario de prueba clave}
Las tareas de ciberseguridad tienen de manera natural una alta densidad de retroalimentación (éxito o fracaso claramente definidos) y un alto riesgo operativo (el costo de una operación errónea es alto).  
Este tipo de tareas permite examinar al mismo tiempo la capacidad de ejecución del agente, su capacidad de seguir restricciones y su capacidad de recuperación.
\end{knowledgebox}

\begin{warningbox}{«Inflar la puntuación» no equivale a una mejora real de la capacidad}
Si la distribución del benchmark es demasiado estrecha, el agente puede aprender rutinas específicas en lugar de una estrategia general.  
El expositor enfatiza que la distribución de tareas y las dimensiones de evaluación deben ampliarse de manera continua, para evitar un progreso aparente de «puntuación alta pero no transferible».
\end{warningbox}

\subsection{Resumen de la sección}
El avance clave de los problem-solving agent proviene de ser «ejecutable + evaluable + revisable». Lo que trabajos del tipo MLAgentBench impulsan realmente es un paradigma de investigación, y no solo el resultado de una sola tabla de clasificación.

